%!TEX program = xelatex
\documentclass[preview]{standalone}
\usepackage{amssymb,amsmath,amsfonts,amsthm,amstext,pxfonts}
\usepackage{graphicx}
\usepackage[usenames,dvipsnames]{xcolor}
\usepackage{subfigure}
\usepackage{tikz,tikz-3dplot,tkz-euclide,circuitikz,siunitx,pstricks-add,pst-coil,pst-3dplot}
\usepackage{pst-plot,pst-func,pst-eucl,pst-solides3d}
\usetikzlibrary{scopes}
\usetikzlibrary{through}
\usetikzlibrary{lindenmayersystems}
\usetikzlibrary[shadings]
\usetikzlibrary{arrows}
\usetikzlibrary{intersections,positioning}

\begin{document}
    \pagenumbering{gobble}
    \newrgbcolor{zzttqq}{0.6 0.2 0}

    % ADICIONADO: 'opacity=1' no psset ativa a chave de transparência no xkeyval do PSTricks
    \psset{xunit=1.0cm,yunit=1.0cm,algebraic=true,dotstyle=o,dotsize=3pt 0,linewidth=0.8pt,arrowsize=3pt 2,arrowinset=0.25,opacity=1}

    \begin{pspicture*}(1,-1.96)(12.08,5.79)

        % PLANO 1 - Muito transparente (10% de cor) utilizando a propriedade 'opacity' correta
        \pspolygon[linecolor=zzttqq,fillcolor=zzttqq,fillstyle=solid,opacity=0.1](2,4)(3.52,4.72)(8.74,-0.54)(7.22,-1.26)

        % PLANO 2 - Transparência média (30% de cor)
        \pspolygon[linecolor=zzttqq,fillcolor=zzttqq,fillstyle=solid,opacity=0.3](4.86,-1.61)(11.05,4.02)(10.05,4.54)(3.86,-1.09)

        % PLANO 3 - Mais visível (60% de cor)
        \pspolygon[linecolor=zzttqq,fillcolor=zzttqq,fillstyle=solid,opacity=0.6](2.32,3.38)(1.42,2.7)(11.54,3.07)(10.78,3.47)

        % Linhas de contorno e interseções
        \psline[linecolor=zzttqq](2,4)(3.52,4.72)
        \psline[linecolor=zzttqq](3.52,4.72)(8.74,-0.54)
        \psline[linecolor=zzttqq](8.74,-0.54)(7.22,-1.26)
        \psline[linecolor=zzttqq](7.22,-1.26)(2,4)
        \psline[linecolor=zzttqq](4.86,-1.61)(11.05,4.02)
        \psline[linecolor=zzttqq](11.05,4.02)(10.05,4.54)
        \psline[linecolor=zzttqq](10.05,4.54)(3.86,-1.09)
        \psline[linecolor=zzttqq](3.86,-1.09)(4.86,-1.61)
        \psline(6.28,-0.32)(6.71,1.5)
        \psline[linecolor=zzttqq](2.32,3.38)(1.42,2.7)
        \psline[linecolor=zzttqq](1.42,2.7)(11.54,3.07)
        \psline[linecolor=zzttqq](11.54,3.07)(10.78,3.47)
        \psline[linecolor=zzttqq](10.78,3.47)(2.32,3.38)
        \psline(4.82,3.41)(3.22,2.77)
        \psline(8.86,3.45)(9.95,3.01)

        \rput[tl](10.45,3.3){$\pi_1$}
        \rput[tl](9.83,4.3){$\pi_2$}
        \rput[tl](3.14,4.47){$\pi_3$}
    \end{pspicture*}
\end{document}
